\section{Method}
\subsection{Amodal candidate detection}
The provided annotations mark the complete object extent, including occluded portions. A pretrained YOLO26n detector was fine-tuned for 80 epochs on the 1,000 labeled training scene pairs as a single object class. Training used 640-pixel input, batch size 16, seed 42, and the supplied training augmentation. The best detector checkpoint was selected by the Ultralytics validation mAP criterion. At inference, the same detector is applied independently to both views at image size 960, confidence 0.45, NMS IoU 0.6, and at most 40 boxes per image. Predicted coordinates are mapped back to the original 512 by 512 image grid. No test labels or manual test annotations are used.

\subsection{Pair features and scene-level matching}
For every candidate box in A and B, we compute ten pair features: the normalized A-to-B and B-to-A camera reprojection residuals, their disagreement, HSV-histogram cosine similarity, DINOv2 ViT-S/14 crop-feature cosine similarity, the two detector confidence scores, and log ratios for width, height, and area. DINOv2 is frozen and used only to extract crop descriptors; it is not fine-tuned \cite{dinov2}. Each box bottom-center is back-projected to the supplied orthographic table plane and projected into the opposite view. This cue is a soft location prior: amodal box geometry and object height mean it is not a reliable pixel correspondence by itself.

The scene matcher encodes each candidate pair, then adds row and column context using the mean and maximum encodings of competing links. It predicts a link score and separate unmatched scores for every box. The utility of linking boxes $i$ and $j$ is $u_{ij}=s_{ij}-d_i^A-d_j^B$, where $s_{ij}$ is the link score and $d_i^A,d_j^B$ are the unmatched scores. A maximum-weight partial one-to-one assignment accepts only links with utility at least $-2.10$; remaining detections become one-sided records. This directly controls whether a box is submitted and prevents a box from appearing in two records, following the assignment structure in the classical assignment problem~\cite{kuhn1955}.

The matcher is a small message-passing network implemented in PyTorch. It is trained with a bidirectional cross-entropy objective: each A box chooses a B box or an unmatched option, and each B box makes the corresponding choice. Training scenes are split deterministically into 800 matcher-fitting scenes and 200 threshold-calibration scenes. Feature normalization and matcher weights are fit on the 800 scenes; the acceptance threshold is selected on the held-out 200 scenes by association F1, with count MAE and exact-count accuracy as tie-breakers. Validation is reserved for final model comparison.
